\documentclass[11pt,a4paper]{article}
\usepackage[T1]{fontenc}
\usepackage{lmodern}
\usepackage[margin=27mm]{geometry}
\usepackage{amsmath,amssymb,amsthm,mathtools}
\usepackage{microtype}
\usepackage{booktabs,array}
\usepackage{enumitem}
\usepackage[hidelinks]{hyperref}
\setlength{\parskip}{5pt plus 1pt}
\setlength{\parindent}{0pt}
\setlength{\emergencystretch}{3em}
\setlist{itemsep=2pt,topsep=3pt}
\raggedbottom
\newtheorem{theorem}{Theorem}[section]
\newtheorem{lemma}[theorem]{Lemma}
\newtheorem{corollary}[theorem]{Corollary}
\newtheorem{proposition}[theorem]{Proposition}
\theoremstyle{definition}
\newtheorem{definition}[theorem]{Definition}
\newtheorem{example}[theorem]{Example}
\theoremstyle{remark}
\newtheorem{remark}[theorem]{Remark}
\newcommand{\Z}{\mathbb{Z}}
\newcommand{\N}{\mathbb{N}}
\newcommand{\Q}{\mathbb{Q}}
\newcommand{\F}{\mathbb{F}}
\newcommand{\ps}[1]{#1[[x]]}
\newcommand{\coeff}[2]{[x^{#1}]\,#2}
\newcommand{\hl}[1]{\textbf{#1}}
\newcommand{\para}[1]{\par\medskip\noindent\textbf{#1}\enspace}
\hypersetup{pdftitle={When are the coefficients of B(B(B(x))) odd? A proof of two conjectures in OEIS A396093},pdfauthor={Dev D. Goyal},pdfsubject={OEIS A396093}}
\title{When are the coefficients of $B(B(B(x)))$ odd?\\[6pt]\large A proof of two conjectures in OEIS A396093}
\author{Dev D. Goyal}
\date{October 2026}
\begin{document}
\maketitle
\begin{abstract}
Substituting the power series $B(x)=x/(1-x)^2=x+2x^2+3x^3+\cdots$ into itself twice gives the series
$A(x)=B(B(B(x)))=x+6x^2+33x^3+174x^4+\cdots$, whose coefficients form the sequence A396093 of the On-Line Encyclopedia
of Integer Sequences (OEIS). Paul D.~Hanna conjectured two rules for when these coefficients are even. We prove both,
and more: the coefficient of $x^n$ is odd if and only if $n$ leaves remainder $1$, $3$, $7$ or $9$ on division by $10$.
The proof works with coefficients modulo $2$, where the triple substitution reduces to the fraction
$(x+x^3+x^7+x^9)/(1+x^{10})$. The result has also been checked by computer and proved formally in the Lean proof
assistant.
\end{abstract}
\begin{center}\begin{minipage}{0.88\linewidth}\small
\textbf{How this paper was produced.} This paper reports the result of an AI research effort directed by
Dev~D.~Goyal. ChatGPT (OpenAI), using the GPT-6 model with Pro reasoning and working autonomously from prompts by
Dev~D.~Goyal, found the proof, wrote its Lean~4 formalization, and wrote a first draft of the paper.
Dev~D.~Goyal chose the problems and asked for proofs, formal verification and revisions, but did not supply proof
ideas. Claude Opus~5.5, used with Claude Code, then checked the work independently: it rebuilt the formal proofs from
source and confirmed that they rely only on Lean's three standard axioms, recomputed the sequence from its definition
compared it with the published terms and tested every claim numerically, and read every step of the proofs. It also wrote the present
version of this paper.
Dev~D.~Goyal has verified the results and is responsible for them. The paper has not been peer reviewed.
The section ``About this paper'' at the end gives the details.
\end{minipage}\end{center}

\section{Introduction}

\subsection{The sequence}\label{sec:sequence}
The \emph{On-Line Encyclopedia of Integer Sequences} (OEIS)~\cite{oeisdb} is a large public database of integer
sequences. Many
entries also record \emph{conjectures}: patterns that the contributor observed in the data but did not prove.

This paper is about entry A396093~\cite{oeis}, contributed by Paul~D.~Hanna in May 2026. It starts from the series
\[
B(x)=\frac{x}{(1-x)^2}=x+2x^2+3x^3+4x^4+\cdots,
\]
whose coefficient of $x^n$ is $n$. Substituting $B$ into itself twice gives a new series,
\[
A(x)=B\bigl(B(B(x))\bigr)=\sum_{n\ge0}a(n)\,x^n=x+6x^2+33x^3+174x^4+892x^5+\cdots .
\]
(Section~\ref{sec:background} explains what substituting one series into another means.) The coefficients $a(n)$ form
the sequence; its first terms, starting at $n=0$, are 0, 1, 6, 33, 174, 892, 4480, 22149, 108144, 522685, 2505112, 11921919, \dots{} The entry reads:
\begin{quote}\small\ttfamily\raggedright G.f. A(x) = B(B(B(x))), the third iteration of B(x) = x/(1-x)\^{}2.\par\smallskip Conjecture: a(2*n) is even for n >= 1.\par\smallskip Conjecture: a(2*n-1) is even iff n is of the form 5*k-2 (k >= 1) for n >= 1.\par\smallskip\hfill\textrm{--- OEIS A396093, Paul D. Hanna, May 2026}\end{quote}

Here ``G.f.''\ stands for \emph{generating function}, the series whose coefficients are the terms of the sequence.

\subsection{The question}
The two conjectures say when $a(n)$ is even. The first says that every term with an even index is even. The second
says that a term with an odd index $2n-1$ is even exactly when $n$ is one of $3,8,13,18,\dots$, that is, when the index is
one of $5,15,25,35,\dots$. The table shows the parity of the first forty terms ($1$ means odd, $0$ means even).
The indices printed in bold are those that leave remainder $1$, $3$, $7$ or $9$ on division by $10$.
\begin{center}\small\setlength{\tabcolsep}{3.2pt}
\begin{tabular}{>{\raggedright\arraybackslash}p{6.4em}>{\centering\arraybackslash}p{1.15em}>{\centering\arraybackslash}p{1.15em}>{\centering\arraybackslash}p{1.15em}>{\centering\arraybackslash}p{1.15em}>{\centering\arraybackslash}p{1.15em}>{\centering\arraybackslash}p{1.15em}>{\centering\arraybackslash}p{1.15em}>{\centering\arraybackslash}p{1.15em}>{\centering\arraybackslash}p{1.15em}>{\centering\arraybackslash}p{1.15em}>{\centering\arraybackslash}p{1.15em}>{\centering\arraybackslash}p{1.15em}>{\centering\arraybackslash}p{1.15em}>{\centering\arraybackslash}p{1.15em}>{\centering\arraybackslash}p{1.15em}>{\centering\arraybackslash}p{1.15em}>{\centering\arraybackslash}p{1.15em}>{\centering\arraybackslash}p{1.15em}>{\centering\arraybackslash}p{1.15em}>{\centering\arraybackslash}p{1.15em}}
\toprule
$n$ & 0 & \hl{1} & 2 & \hl{3} & 4 & 5 & 6 & \hl{7} & 8 & \hl{9} & 10 & \hl{11} & 12 & \hl{13} & 14 & 15 & 16 & \hl{17} & 18 & \hl{19} \\ \midrule
$a(n)\bmod 2$ & 0 & \hl{1} & 0 & \hl{1} & 0 & 0 & 0 & \hl{1} & 0 & \hl{1} & 0 & \hl{1} & 0 & \hl{1} & 0 & 0 & 0 & \hl{1} & 0 & \hl{1} \\
\bottomrule
\end{tabular}
\\[4pt]
\begin{tabular}{>{\raggedright\arraybackslash}p{6.4em}>{\centering\arraybackslash}p{1.15em}>{\centering\arraybackslash}p{1.15em}>{\centering\arraybackslash}p{1.15em}>{\centering\arraybackslash}p{1.15em}>{\centering\arraybackslash}p{1.15em}>{\centering\arraybackslash}p{1.15em}>{\centering\arraybackslash}p{1.15em}>{\centering\arraybackslash}p{1.15em}>{\centering\arraybackslash}p{1.15em}>{\centering\arraybackslash}p{1.15em}>{\centering\arraybackslash}p{1.15em}>{\centering\arraybackslash}p{1.15em}>{\centering\arraybackslash}p{1.15em}>{\centering\arraybackslash}p{1.15em}>{\centering\arraybackslash}p{1.15em}>{\centering\arraybackslash}p{1.15em}>{\centering\arraybackslash}p{1.15em}>{\centering\arraybackslash}p{1.15em}>{\centering\arraybackslash}p{1.15em}>{\centering\arraybackslash}p{1.15em}}
\toprule
$n$ & 20 & \hl{21} & 22 & \hl{23} & 24 & 25 & 26 & \hl{27} & 28 & \hl{29} & 30 & \hl{31} & 32 & \hl{33} & 34 & 35 & 36 & \hl{37} & 38 & \hl{39} \\ \midrule
$a(n)\bmod 2$ & 0 & \hl{1} & 0 & \hl{1} & 0 & 0 & 0 & \hl{1} & 0 & \hl{1} & 0 & \hl{1} & 0 & \hl{1} & 0 & 0 & 0 & \hl{1} & 0 & \hl{1} \\
\bottomrule
\end{tabular}
\end{center}

The data suggest a single rule for all indices, which we prove.

\subsection{What this paper proves}
\textbf{Main result} (Theorem~\ref{thm:main}). \emph{The coefficient $a(n)$ is odd if and only if $n$ leaves remainder
$1$, $3$, $7$ or $9$ on division by $10$.}

Both of Hanna's conjectures follow at once (Corollary~\ref{cor:conj}). The key step is a short calculation showing that,
once every coefficient is replaced by its remainder modulo $2$, the series $A(x)$ becomes the fraction
\[
\frac{x+x^3+x^7+x^9}{1+x^{10}}=(x+x^3+x^7+x^9)\,(1+x^{10}+x^{20}+\cdots),
\]
from which the pattern can be read off directly.

\subsection{What was known}
The OEIS entry also records a closed formula for $A(x)$ as a ratio of two polynomials (due to Vaclav Kotesovec), and a
linear recurrence for $a(n)$ (due to Seiichi Manyama). For any sequence given by such a ratio, the remainders modulo
$2$ are known to repeat periodically from some point on: there are only finitely many possible ``states'' of the
recurrence modulo $2$. So a periodic parity pattern was to be expected, and in principle its period could be found by a
finite computation from the recorded formula. We do not rely on that formula. The proof below works directly from the
definition and shows \emph{why} the period is $10$. We are not aware of an earlier proof of the two conjectures.

\subsection{Motivation}
Substituting a series into itself is a basic operation: it describes what happens when a process is repeated, and
iterates of simple series like $B$ appear throughout combinatorics. Their coefficients quickly become large and
irregular. Remainders modulo small numbers are often the first place where structure becomes visible, and explaining
such patterns is a standard way of understanding a sequence. The example also shows a general technique: computing
modulo $2$ can collapse a complicated expression (here a ratio of two polynomials of degrees $7$ and $8$) into something
transparent. Section~\ref{sec:discussion} describes related questions.

\subsection{How to read this paper}
Section~\ref{sec:background} explains the notions we need: power series, substitution, and remainders. Section~\ref{sec:results}
states the results precisely, and Section~\ref{sec:proofs} proves them. Section~\ref{sec:checks} describes how the results were
checked, by computer and by a formal proof in Lean. Section~\ref{sec:discussion} discusses the scope of the result and
open questions.
\section{Background}\label{sec:background}
This section collects the definitions and facts used later. Readers who know formal power series and congruences can skip it and return when needed.
\subsection{Formal power series}
A \emph{formal power series} is an expression
\[
U(x)=u_0+u_1x+u_2x^2+u_3x^3+\cdots ,
\]
determined by its sequence of coefficients $u_0,u_1,u_2,\dots$; see~\cite{niven,stanley}. No number is ever
substituted for $x$, so questions of convergence do not arise. We write $\coeff{n}{U}=u_n$ for the coefficient of $x^n$ in $U$; two series are equal
when all their coefficients are equal.

Series are added coefficient by coefficient and multiplied like polynomials:
\[
\coeff{n}{(UV)}=u_0v_n+u_1v_{n-1}+\cdots+u_nv_0 .
\]
Each coefficient of a product is a finite sum. For example, $(1-x)(1+x+x^2+x^3+\cdots)=1$, so $1+x+x^2+\cdots$ is
the \emph{multiplicative inverse} of $1-x$, written $\frac{1}{1-x}$. More generally, every series $D$ with constant term
$1$ has a unique multiplicative inverse $\frac1D$, whose coefficients are determined one at a time by $D\cdot\frac1D=1$
using only addition, subtraction and multiplication. Whenever a quotient of power series appears in this paper, its
denominator has constant term $1$, so the quotient is again a power series.

A \emph{commutative ring} is a set with an addition, a subtraction and a multiplication that satisfy the usual rules
of arithmetic; examples are $\Z$, $\Q$ and the integers modulo $m$. For a commutative ring $R$ we write $\ps{R}$ for
the set of power series with coefficients in $R$; for instance, writing $\F_2$ for the integers modulo $2$, $\ps{\F_2}$
consists of the series with coefficients modulo $2$. We write
$U(0)$ for the constant term $u_0$ of $U$, and $U=O(x^d)$ when the coefficients of $1,x,\dots,x^{d-1}$ in $U$ are all
zero. Thus $U-V=O(x^d)$ means that $U$ and $V$ have the same coefficient of $x^n$ for every $n<d$. A series with
constant term $0$ and coefficient of $x$ equal to $1$, that is, a series of the form $x+u_2x^2+u_3x^3+\cdots$, is
called \emph{normalized}.

\subsection{Substituting one series into another}
If $V(x)$ has constant term $0$, we can substitute $V$ for $x$ in $U$:
\[
U(V(x))=u_0+u_1V(x)+u_2V(x)^2+u_3V(x)^3+\cdots .
\]
This is called \emph{composition}, also written $U\circ V$. It makes sense because $V(x)^k$ has no terms below $x^k$,
so only the finitely many terms with $k\le n$ affect the coefficient of $x^n$. For example, if $U=x+x^2$ and $V=x^2$, then
$U(V(x))=x^2+x^4$. Composition is associative: $U\circ(V\circ W)=(U\circ V)\circ W$.

Substituting a series into itself repeatedly gives its \emph{iterates}: $U^{\circ2}(x)=U(U(x))$,
$U^{\circ3}(x)=U(U(U(x)))$, and so on, with $U^{\circ0}(x)=x$ and $U^{\circ1}=U$. They must not be confused with
\emph{powers} $U^2=U\cdot U$, $U^3=U\cdot U\cdot U$. (The OEIS sometimes writes $A^k(x)$ for the $k$-th iterate. In
this paper the small circle always means an iterate, and a plain exponent always means a power.)

\subsection{Remainders and congruences}
For an integer $m\ge2$ we write $a\equiv b\pmod m$ (``$a$ is congruent to $b$ modulo $m$'') when $a-b$ is divisible by
$m$, that is, when $a$ and $b$ leave the same remainder after division by $m$. The remainder of $a$ on division by
$m$ is called its \emph{residue} modulo $m$; it is one of $0,1,\dots,m-1$.

Addition and multiplication are compatible with congruences, so the residues $0,1,\dots,m-1$, added and multiplied
modulo $m$, form a ring, written $\Z/m\Z$. When $m=p$ is a prime we also write $\F_p$. In $\F_2=\{0,1\}$ we have $1+1=0$,
so addition and subtraction coincide.

For power series with integer coefficients, $U\equiv V\pmod m$ means that \emph{every} coefficient of $U-V$ is
divisible by $m$. Equivalently, the two series become equal when every coefficient is replaced by its residue; we call
this \emph{reducing} the series modulo $m$. The key fact is the following.

\emph{Reduction modulo $m$ is compatible with all the operations above.} To reduce a sum, a product, a substitution, or
an inverse, we may reduce the ingredients first and then compute in $\Z/m\Z$. This holds because every
coefficient of the result is obtained from finitely many coefficients of the ingredients by adding, subtracting and
multiplying, and these operations are compatible with remainders.

\subsection{The Frobenius identity}
\begin{lemma}[Frobenius identity]\label{lem:frobenius}
Let $p$ be a prime. For every power series $U$ with coefficients in $\F_p$, we have $U(x)^p=U(x^p)$.
\end{lemma}

\begin{proof}
In $\F_p$, the binomial theorem simplifies to $(a+b)^p=a^p+b^p$, because every binomial coefficient
$\binom pi$ with $0<i<p$ is divisible by $p$. Also $c^p=c$ for every $c\in\F_p$ (Fermat's little theorem). So
$\bigl(\sum_nu_nx^n\bigr)^p=\sum_nu_n^px^{pn}=\sum_nu_nx^{pn}=U(x^p)$. The first equality holds for polynomials by
repeated use of $(a+b)^p=a^p+b^p$; it holds for power series because each coefficient of $U(x)^p$ involves only finitely
many coefficients of $U$.
\end{proof}

For $p=2$ this says that squaring a series modulo $2$ doubles every exponent: for instance,
$(x+x^3)^2\equiv x^2+x^6\pmod2$. In particular, the square of any series has, modulo $2$, no terms of odd degree.

\section{Results}\label{sec:results}

\begin{theorem}\label{thm:main}
For every $n\ge0$,
\[
a(n)\ \text{is odd}\quad\Longleftrightarrow\quad n\equiv 1,\ 3,\ 7\ \text{or}\ 9 \pmod{10}.
\]
Equivalently, modulo $2$ the generating function is
\begin{equation}\label{eq:key}
A(x)\equiv\frac{x+x^3+x^7+x^9}{1+x^{10}}=\sum_{j\ge0}\bigl(x^{10j+1}+x^{10j+3}+x^{10j+7}+x^{10j+9}\bigr)\pmod 2 .
\end{equation}
\end{theorem}

\begin{corollary}[Hanna's conjectures]\label{cor:conj}
\begin{enumerate}[label=(\alph*)]
\item $a(2n)$ is even for every $n\ge0$.
\item For $n\ge1$, $a(2n-1)$ is even if and only if $n=5k-2$ for some integer $k\ge1$.
\end{enumerate}
\end{corollary}

\begin{proof}
(a) An even number never leaves remainder $1$, $3$, $7$ or $9$ on division by $10$.
(b) An odd number $2n-1$ leaves one of the remainders $1,3,5,7,9$. By Theorem~\ref{thm:main}, $a(2n-1)$ is even
if and only if $2n-1\equiv5\pmod{10}$. This means $2n\equiv6\pmod{10}$, that is, $n\equiv3\pmod5$. For
$n\ge1$ these are the numbers $3,8,13,\dots$, which are exactly the numbers $5k-2$ with $k\ge1$.
\end{proof}

\section{Proofs}\label{sec:proofs}

\para{The idea in brief.} We reduce all coefficients modulo $2$. There, $B(x)=x/(1+x^2)$, and substituting a fraction
into $B$ gives another fraction by a simple rule (Lemma~\ref{lem:step}). Applying the rule three times expresses
$A=B\circ B\circ B$ as a fraction, and multiplying its numerator and denominator by $1+x^2$ turns the denominator into
$1+x^{10}$.

From now on we work in $\ps{\F_2}$, that is, with coefficients modulo $2$, and we keep the names $A$ and $B$ for the
reduced series. Because reduction is compatible with substitution, the reduction of $A=B(B(B(x)))$ equals the triple
substitution of the reduction of $B$. Since $-1=1$ in $\F_2$, and since $(1+x)^2=1+x^2$ there (Lemma~\ref{lem:frobenius}),
\[
B(x)=\frac{x}{(1-x)^2}=\frac{x}{(1+x)^2}=\frac{x}{1+x^2}.
\]

\begin{lemma}\label{lem:step}
Let $P$ and $D$ be series over $\F_2$ with $P(0)=0$ and $D(0)=1$. Then
\[
B\!\left(\frac PD\right)=\frac{P\,D}{D^2+P^2}.
\]
\end{lemma}

\begin{proof}
Put $T=P/D$. Then $B(T)=T/(1+T^2)$. Multiplying the numerator and the denominator by $D^2$ gives
$PD/(D^2+P^2)$. The new denominator has constant term $1$, so this fraction is again a power series.
\end{proof}

\begin{proof}[Proof of Theorem~\ref{thm:main}]
Starting from $x=x/1$ and applying Lemma~\ref{lem:step} three times (all arithmetic modulo $2$):
\begin{center}\small\setlength{\tabcolsep}{4pt}\renewcommand{\arraystretch}{1.2}
\begin{tabular}{lll}
\toprule
series & numerator $P$ & denominator $D$\\
\midrule
$B(x)$ & $x$ & $1+x^2$\\
$B(B(x))$ & $x(1+x^2)=x+x^3$ & $(1+x^2)^2+x^2=1+x^2+x^4$\\
$B(B(B(x)))$ & $(x+x^3)(1+x^2+x^4)=x+x^7$ & $(1+x^2+x^4)^2+(x+x^3)^2=1+x^2+x^4+x^6+x^8$\\
\bottomrule
\end{tabular}
\end{center}
For instance, in the last row $(x+x^3)(1+x^2+x^4)=x+2x^3+2x^5+x^7$, and the terms with coefficient $2$ vanish
modulo $2$. Hence
\[
A(x)=\frac{x+x^7}{1+x^2+x^4+x^6+x^8}.
\]
Now multiply the numerator and the denominator by $1+x^2$. The numerator becomes $(x+x^7)(1+x^2)=x+x^3+x^7+x^9$. The
denominator becomes $(1+x^2)(1+x^2+x^4+x^6+x^8)=1+x^{10}$, because every middle term appears twice and so vanishes.
This proves the first form of \eqref{eq:key}. Finally, $1/(1+x^{10})=1+x^{10}+x^{20}+\cdots$, which
gives the second form. Reading off coefficients, $a(n)$ is odd exactly when $n\equiv1,3,7,9\pmod{10}$.
\end{proof}

\section{Checking the results}\label{sec:checks}

A mathematical proof can contain mistakes, so the results were checked in two independent ways.

\para{Numerical check.}
We computed $B(B(B(x)))$ directly from the definition. The result reproduces all 27 terms printed in the OEIS
entry, and the parity rule of Theorem~\ref{thm:main} holds for every $n<600$.
The computation was done with the PARI/GP computer algebra system, using code that was written separately from the
formal proofs (file \texttt{verification/independent\_checks.gp} in the accompanying repository~\cite{repo}).
A computation can only test finitely many cases, so it is a check on the definitions and on the indexing, not a proof.

\para{Formal proof in Lean.}
Lean~\cite{lean} is a \emph{proof assistant}: a program in which mathematical statements and their proofs are written
in a precise formal language, and every step is checked mechanically down to the basic rules of logic. The theorems of Section~\ref{sec:results} are proved in Lean~4 (version 4.19.0) on top of
Mathlib~\cite{mathlib}, Lean's library of formalized mathematics. The statements are collected in the file \texttt{lean/Entries/A396093.lean} of the repository,
in the namespace \texttt{OEIS.A396093}. There $B$ is defined by the equation $B\cdot(1-x)^2=x$ and $A$ as $B\circ B\circ B$.
The table lists where each result of this paper is proved.
\begin{center}\small\renewcommand{\arraystretch}{1.25}
\begin{tabular}{>{\raggedright\arraybackslash}p{0.52\linewidth}>{\raggedright\arraybackslash}p{0.4\linewidth}}
\toprule
Statement in this paper & Lean theorem(s)\\
\midrule
$B(x)\cdot(1-x)^2=x$ (definition of $B$) & \texttt{source\_\allowbreak{}denominator} \\
$A=B\circ B\circ B$, and $a(n)$ is its coefficient of $x^n$ & \texttt{source\_\allowbreak{}composition}, \texttt{coefficient\_\allowbreak{}definition} \\
Equation~\eqref{eq:key}: $A\cdot(1+x^{10})=x+x^3+x^7+x^9$ over $\F_2$ & \texttt{reduced\_\allowbreak{}denominator} \\
Theorem~\ref{thm:main} & \texttt{complete\_\allowbreak{}parity}, \texttt{even\_\allowbreak{}iff} \\
Corollary~\ref{cor:conj}(a) and (b) & \texttt{even\_\allowbreak{}indices}, \texttt{odd\_\allowbreak{}indices} \\
\bottomrule
\end{tabular}
\end{center}
The Lean project was rebuilt from source, and the axioms on which each of these theorems depends were listed with
Lean's \texttt{\#print axioms} command. Every one depends only on the three standard axioms of Lean's logic
(\texttt{propext}, \texttt{Classical.choice}, \texttt{Quot.sound}); no step is left unproved. A formal proof certifies
exactly the statement that Lean checks. Whether that statement says what the OEIS entry means is a separate question,
so the Lean definitions were also compared with the definition in the OEIS entry
(Section~\ref{sec:sequence}).

\section{Discussion}\label{sec:discussion}

\para{Scope.} Theorem~\ref{thm:main} describes the parity of every coefficient, so it settles both conjectures of the
entry completely. It says nothing about the remainders of $a(n)$ modulo other numbers, or about the size of $a(n)$.

\para{Other iterates.} The same calculation applies to every iterate of $B$. The first two rows of the table above
already show that $B(x)\equiv x/(1+x^2)$ and, after multiplying by $1+x^2$, that
$B(B(x))\equiv(x+x^5)/(1+x^6)\pmod2$. So the coefficients of $B(B(x))$ (OEIS A030267) are odd exactly when
$n\equiv1$ or $5\pmod 6$. For the fourth iterate, a computer calculation shows a parity pattern of period $34$. We do not
know a rule that gives the period for every iterate.

\para{Higher moduli.} In a computer calculation, the remainders of $a(n)$ modulo $4$ repeat with period $20$,
starting from $n=0$; we have not proved that this continues. Since $A$ is a ratio of polynomials, these remainders are
periodic from some point on, and the method of this paper could be used to determine the period.

\section*{About this paper}\label{sec:about}
\para{How the work was done.}
This paper is one of 29 papers on conjectures from the OEIS that came out of an AI research effort directed by
Dev~D.~Goyal~\cite{repo}. The roles were as follows.
\begin{itemize}
\item \emph{Dev~D.~Goyal} started the project, chose which problems to pursue, asked for proofs, formal verification and
revisions, and decided what to publish. Dev~D.~Goyal did not supply proof strategies or intermediate arguments.
\item \emph{ChatGPT} (OpenAI), using the GPT-6 model with Pro reasoning, worked autonomously from those prompts. It
searched the OEIS for suitable conjectures, found the proofs, wrote the Lean~4 formalization, and wrote a first draft
of this paper.
\item \emph{Claude Opus~5.5}, used with Claude Code, checked the work independently. It rebuilt the Lean project from source and confirmed that the theorems listed in Section~\ref{sec:checks} rely only on
Lean's standard axioms; compared the Lean definitions with the OEIS definition; recomputed the sequence with separately
written PARI/GP code, compared it with the terms published in the OEIS, and tested every claim of the paper on many
terms; and read every step of the proofs. It then wrote the present version of the paper, expanding the
explanations and, in places, writing out an argument in more detail or by a more elementary route. The theorems of
Section~\ref{sec:results} and the proof strategies are ChatGPT's. The Discussion section, including the short arguments
in it, was written by Claude Opus~5.5; its observations were checked by computer and are not part of the Lean
formalization.
\end{itemize}
Dev~D.~Goyal has verified the results and is responsible for this paper.

\para{Review status.} This paper has not been peer reviewed, and it has not been submitted to a journal or to the OEIS.

\para{Priority.} An OEIS entry that labels a statement as a conjecture shows only that its contributor did not know
a proof; it does not show that no proof existed. During the research only limited searches for earlier work were
made. No claim of priority is made.

\para{Corrections.} Please report any error in the issue tracker of the repository~\cite{repo}, where the Lean sources and the numerical checks are also
available.

\section*{Acknowledgements}
The sequence and both conjectures are due to Paul~D.~Hanna~\cite{oeis}. The closed formula and recurrence mentioned in
the introduction were contributed to the entry by Vaclav~Kotesovec and Seiichi~Manyama.

\begin{thebibliography}{99}\small
\raggedright
\bibitem{oeis} P.~D.~Hanna, entry A396093 (Third iteration of $x/(1-x)^2$), \emph{The On-Line Encyclopedia of Integer Sequences}, May 2026. \url{https://oeis.org/A396093}. Text as of October 2026.
\bibitem{oeisdb} N.~J.~A.~Sloane and The OEIS Foundation Inc., \emph{The On-Line Encyclopedia of Integer Sequences}. \url{https://oeis.org}.
\bibitem{stanley} R.~P.~Stanley, \emph{Enumerative Combinatorics}, Volume~2, Cambridge University Press, 1999. (Chapter~5 treats formal power series, composition and compositional inverses.)
\bibitem{niven} I.~Niven, Formal power series, \emph{American Mathematical Monthly} 76 (1969), 871--889.
\bibitem{lean} L.~de~Moura and S.~Ullrich, The Lean~4 theorem prover and programming language, in \emph{Automated Deduction -- CADE 28}, Lecture Notes in Computer Science 12699, Springer, 2021, 625--635.
\bibitem{mathlib} The mathlib Community, The Lean mathematical library, in \emph{Proceedings of the 9th ACM SIGPLAN International Conference on Certified Programs and Proofs (CPP 2020)}, 367--381. Version used: commit \texttt{c44e0c8ee63c}.
\bibitem{repo} D.~D.~Goyal, \emph{Proofs and corrections of OEIS conjectures}, repository, 2026. \url{https://github.com/dvgyl/oeis-a396843-proof}.
\end{thebibliography}
\end{document}
